% Part I: a first tour with the demo. Sources: docs/getting-started.md; docs/administration.md
% ("The demo"); tools/make-demo-data.py (what the demo holds); medorama/src/main.rs (init's and
% demo add's messages), demo.rs.
\chapter{Quick start: the demo}
\label{chap:quickstart}

This chapter takes a first-time user or an administrator from a fresh clone of the repository
to a running Medorama with invented patients, and then through a short tour: a patient, a
cohort, a recording and Privacy Mode. It ends by removing the demo. The whole procedure takes
about fifteen minutes, most of which is the build. Deploying Medorama for other people is the
subject of \cref{chap:administration} and \cref{chap:deployment}.

\begin{note}[Research use only]
Medorama is research software, not a medical device. Try it on the demo first, and put real
patient data in it only on a server that your institution has approved (\cref{chap:security}).
\end{note}

\section{What the demo holds}
\label{sec:quickstart-demo}

The \emph{demo} is a set of invented patients for trying Medorama and for showing what it does.
Everything in it is synthetic: \file{tools/make-demo-data.py} invents the names, numbers,
signals and pictures. The demo holds:
\begin{itemize}
\item six patients, \code{P900001} to \code{P900006} (MRNs \code{SYN0000001} to
  \code{SYN0000006}), with their encounters and bed stays, laboratory results, vital signs,
  medications, diagnoses, microbiology, allergies and notes;
\item short ICU recordings (WFDB collections with ECG, arterial pressure, SpO2 pleth,
  respiration, vital signs, alarms and R-wave markers) and a 24-hour Holter ECG;
\item an overnight sleep study for patient \code{P900004}: a two-hour EDF+ file
  (\file{sleep/P900004\_psg\_20240305.edf}) with EEG, EOG, ECG, airflow, thoracic effort and
  SpO2 signals, scored with sleep stages, obstructive apneas, oxygen desaturations and arousals
  as annotations. The study has its own encounter in the sleep laboratory (\code{E9100010}), its
  recording starts at 22:30 New York time, and a note, the \emph{Sleep study report}, sums it up;
\item 12-lead ECGs (WFDB, GE MUSE XML, DICOM, PDF and PNG), chest X-rays and CT series, echo
  loops and reports, and a coronary angiography run;
\item another hospital's OMOP database, which adds earlier stays and results to two of the
  patients.
\end{itemize}

In the catalog, the demo is a cohort named \code{DEMO} that holds its patients. Every user sees
this cohort, read-only, and each demo patient's page is marked \ui{DEMO}. The demo's tables are
imported under the source \code{DEMO}, one database per table (\code{patients},
\code{encounters}, \code{labs}, \code{vitals}, \code{ecg\_index}, \code{icu\_index},
\code{holter\_index}, \code{sleep\_index} and the others), and its OMOP database is connected
as \code{demo\_omop}.

\begin{note}
The demo is off by default: a catalog starts empty unless the demo was asked for. Install it
only where Medorama is tried, demonstrated or taught, never on a server meant for real patients.
\end{note}

\section{Install with the demo}
\label{sec:quickstart-install}

Medorama needs Git, Python 3.9 or later, a C compiler and Rust 1.88 or later. The installer
installs Rust in the user's account when it is missing (\cref{chap:installation} gives the
commands for each system, and the website gives them at
\url{https://medorama.org/docs/install.html}). It builds Medorama, one VitalStream waveform viewer and the database
tools into a \file{.venv} folder inside the clone.

On Linux or macOS, the following commands clone the repository, run the installer with the Rust
viewer and the demo (this takes a few minutes), and activate the environment:
\begin{lstlisting}[style=shell]
git clone https://github.com/alphanumericslab/medorama.git && cd medorama
tools/install-medorama.sh --viewer rs --demo
source .venv/bin/activate
\end{lstlisting}
Activate the environment again in every new terminal.

On Windows, the same steps run in PowerShell, with the Microsoft C++ Build Tools installed:
\begin{lstlisting}[language=PowerShell]
git clone https://github.com/alphanumericslab/medorama.git; cd medorama
powershell -ExecutionPolicy Bypass -File tools\install-medorama.ps1 -Viewer rs -Demo
.venv\Scripts\Activate.ps1
\end{lstlisting}

The option \opt{-{}-viewer rs} builds the Rust VitalStream viewer, which needs nothing beyond
Rust. The default, the C++ viewer, also needs a C++20 compiler. With \opt{-{}-demo}, the
installer writes the demo's files (about 60\,MB) into \file{.venv/share/medorama/demo} and names
that folder in \file{.venv/bin/medorama.demo}, where \cmd{medorama init} finds it.

\section{Create the catalog}
\label{sec:quickstart-init}

The \emph{catalog} is Medorama's own database: its patients, their encounters and data, the
links to their files, and the users. Its settings and database live in a \emph{home folder}, by
default \file{\textasciitilde/.local/share/medorama} (on Windows
\file{\%LOCALAPPDATA\%\textbackslash medorama}). A catalog also needs a \emph{results folder},
which holds the recordings' indexes and the results of actions.

The following command creates the catalog with its results folder:
\begin{lstlisting}[style=shell]
medorama init --results ~/medorama-results
\end{lstlisting}
On Windows, the same command in PowerShell is:
\begin{lstlisting}[language=PowerShell]
medorama init --results $HOME\medorama-results
\end{lstlisting}

\cmd{init} asks twice for the password of the first user, \code{admin}, who has the PHI
permission. The password needs at least 10 characters. \cmd{init} then says where the catalog
is and, because the installer made the demo, loads it as that administrator:
\begin{lstlisting}[style=text]
Catalog created in /home/user/.local/share/medorama
  settings: /home/user/.local/share/medorama/medorama.toml
  database: ... (holds identifiers: keep it private)
  results:  /home/user/medorama-results
  indexes:  /home/user/medorama-results/index
  admin:    admin (with PHI permission)
Later commands use this catalog without --home (remembered in ...).
Loading the demo from .../.venv/share/medorama/demo ...
Demo loaded: 6 patient(s), ... encounter(s), ... item(s), ... recording(s) from ... table(s),
and the OMOP database demo_omop; the cohort DEMO holds them
Deleting the cohort DEMO (an admin) removes all of it: medorama demo remove does the same.
\end{lstlisting}
The other options of \cmd{init} (\opt{-{}-home}, \opt{-{}-index-dir}, \opt{-{}-timezone},
\opt{-{}-port}, a PostgreSQL database and \opt{-{}-no-demo}) are described in
\cref{chap:administration}.

\section{Serve, and sign in}
\label{sec:quickstart-serve}

The following command starts the server, which serves the web page and the API and indexes the
demo's recordings in the background:
\begin{lstlisting}[style=shell]
medorama serve
\end{lstlisting}
\cmd{serve} prints the address of the page, the catalog and the viewer it runs:
\begin{lstlisting}[style=text]
Medorama: open http://127.0.0.1:8090/
  catalog: /home/user/.local/share/medorama
  database: ...
  viewer: rs (started on first use, in the VitalStream tab)
  ...
\end{lstlisting}
The server listens only on this computer (127.0.0.1) by default. To reach it from another
computer, use an SSH tunnel or an HTTPS reverse proxy (\cref{chap:deployment}).
\key{Ctrl}+\key{C} in the terminal stops the server.

To sign in:
\begin{enumerate}
\item Open \url{http://127.0.0.1:8090} in a browser on the same computer.
\item Type \code{admin} as the \ui{User name}, and the password chosen above.
\item Click \ui{Sign in}.
\end{enumerate}
The page opens on the \ui{Cohorts} tab. The administrator holds the \emph{PHI permission}, the
right to see identifiers such as names, medical record numbers (MRNs) and dates of birth. For
administrators, \emph{Privacy Mode} starts off, so the demo's invented names are visible. For
every other user, Privacy Mode starts on (\cref{chap:protecting}).

The top bar has three ways to the patients (\ui{Cohorts}, \ui{Patient Catalog},
\ui{Databases}), the viewers' tabs (\ui{ECG}, \ui{PACS}, \ui{VitalStream}), and the work tabs at
the right (\ui{History}, \ui{Actions}, \ui{Jobs}, \ui{Review}, \ui{Import}, \ui{Admin}). The
\ui{Privacy Mode} switch and the theme button are at the far right (\cref{chap:page}).

While the server runs, its workers index the demo's recordings within seconds. The recordings'
overviews, and the heart rate that the index job derives from the Holter's ECG, appear once that
is done. The \ui{Jobs} tab shows the queue (\cref{chap:actions}).

\section{A first tour}
\label{sec:quickstart-tour}

\subsection{A patient}
\label{sec:quickstart-patient}

To open a patient's page, click the cohort \ui{DEMO} in the list on the left of \ui{Cohorts},
then a patient among its members, such as \ui{Doe, Jane} (patient ID \code{P900001}). You can
also open \ui{Patient Catalog}, which lists every patient with the kinds of data each one has,
and click a row. The patient's page has the following parts (\cref{chap:patients}):
\begin{itemize}
\item The \emph{patient bar} at the top shows the name, sex, age and date of birth, the patient
  ID and pseudonym, the MRN, the span of the data, any allergies, and the mark \ui{DEMO}.
\item The patient column on the left lists the encounters, the medications of the latest stay
  and the conditions.
\item The \emph{Patient Journey} in the middle draws every item on one time line, newest on
  top, with one column per clinical category (ICU, Vitals, Labs, ECG, Holter, Sleep and so on).
  Hover over a mark to see what it is, and click it to open it.
\item On the right, the \emph{report} lists every entry by day, newest first, and the
  \emph{detail} shows the entry chosen.
\end{itemize}

Try the following on patient \code{P900001}:
\begin{enumerate}
\item Click a \emph{12-lead ECG} entry in the report. The detail shows its measurements, the
  interpretation and the tracing on paper. Its button \ui{Open in ECG} opens it in the \ui{ECG}
  tab, with layouts, gain and calipers.
\item Type \code{creatinine} in the report's search box. The report keeps the entries that
  mention it, with the words highlighted.
\item Click an \emph{ICU monitoring} entry. The detail shows the recording's main vital signs
  over its whole length, and a click on them opens VitalStream at that moment
  (\cref{sec:quickstart-recording}).
\item Click an encounter on the left. The journey then shows only that encounter, unfolded, and
  the encounter's chip above the journey (with $\times$) shows all of them again.
\end{enumerate}

Patient \code{P900004} (Ade Okafor) has a 24-hour Holter, whose entry shows the heart rate found
from its ECG over the whole day, and an overnight EDF+ sleep study in the \ui{Sleep} column of
the journey (\cref{fig:patients-sleep}). The sleep study belongs to its own encounter in the
sleep laboratory, whose \emph{Sleep study report} appears among the notes. Other patients have echo loops, CT series and an
angiography run, which open in the \ui{PACS} tab.

\subsection{A cohort}
\label{sec:quickstart-cohort}

A \emph{cohort} is a named set of patients, found across every database by criteria and kept.
To build one:
\begin{enumerate}
\item Open \ui{Cohorts}. The page opens on a new cohort, with the criteria builder and one
  empty expression line.
\item Type the following expression:
\begin{lstlisting}[style=text]
labs.creatinine > 1.5 OR has icu
\end{lstlisting}
  The expression is checked as it is typed, and its meaning, in words, appears beside it. The
  patients who meet it are listed below, each with the reason it matched.
\item Type a name in the field \ui{New cohort's name}, and click \ui{Save all \dots{} as a new
  cohort}.
\end{enumerate}
The cohort's page then shows its definition, its members, and what can be done with it
(refresh, clone, merge, share, export). The following expressions are other examples to try on
the demo:
\begin{lstlisting}[style=text]
has holter
has sleep
female AND has ecg
in ecg_index AND NOT has icu
\end{lstlisting}
\Cref{chap:cohorts} describes the expression language, every operator and every action on a
cohort.

\subsection{A recording}
\label{sec:quickstart-recording}

From an \emph{ICU monitoring} entry of a patient's page, click the vital signs in the detail at
the moment you want to see, or click \ui{Open in VitalStream}. The \ui{VitalStream} tab opens the
recording at that moment, with every channel, the vital signs, the alarms and the breaks on one
time axis. Zoom with the mouse wheel, from the whole stay down to single samples. The tab's
title says whose recording it is: \ui{$\langle$ Patient} returns to the entry on the patient's
page, and \ui{$\langle$} and \ui{$\rangle$} step to the patient's other recordings
(\cref{chap:viewers-in-medorama}).

Open the sleep study of patient \code{P900004} the same way. Its EDF+ annotations (sleep stages,
apneas, desaturations and arousals) appear as events in VitalStream's \ui{Annotations} lane.
VitalStream's own controls are described in \cref{chap:vitalstream}.

The first time a recording is opened, it may still be indexing. Its status appears in the
\ui{Jobs} tab.

\subsection{Privacy Mode}
\label{sec:quickstart-privacy}

\emph{Privacy Mode} shows the same data without any identifier. The server, not the browser,
removes the identifiers, so none of them reaches the page. As \code{admin}, the session starts
out of Privacy Mode, and identifiers are shown. To see Privacy Mode:
\begin{enumerate}
\item Click \ui{Privacy Mode} at the right of the top bar. The page reloads.
\item Look at Jane Doe's page again.
\end{enumerate}
The following has changed:
\begin{itemize}
\item The name is \code{[Patient Name]}, the MRN \code{[MRN]} and the date of birth
  \code{[Date of Birth]}.
\item Patient and encounter IDs are replaced by \emph{pseudonyms}: random codes such as
  \code{MO-D5QM-Y1X5} for a patient and \code{MO-E-QMC6-E2VA} for an encounter.
\item Every date moves back, by the catalog's offset (about a hundred years by default) plus the
  patient's own number of days, so the intervals between the patient's events stay the same.
\item File names become tokens such as \code{mo-f0000.hea}, and the ECG and PACS tabs receive
  files that the server has de-identified.
\item Staff names in notes become \code{[Staff Name]}, and a search for a name or an MRN is
  refused.
\end{itemize}
Click \ui{Privacy Mode} again to leave it.

\screenshot{privacy.png}{A demo patient's page in Privacy Mode: a pseudonym and placeholders
in place of identifiers, and dates moved back.}{fig:quickstart-privacy}

A user without the PHI permission always works this way. To see it for yourself, add such a user
from a second terminal, with the environment activated. The following command adds a user with
the role \code{viewer} and asks for the new user's password:
\begin{lstlisting}[style=shell]
medorama user add reader --role viewer
\end{lstlisting}
Then sign in as \code{reader} in a private browser window. \Cref{chap:protecting} explains how
the privacy gate works, and what it cannot clean (for example, text burned into images).

\paragraph{The theme.} The theme button at the far right of the top bar switches the page
between its light and its dark theme. Until you choose, the page follows the system's setting,
and the viewers in its tabs follow the page (\cref{chap:page}).

\section{Removing the demo}
\label{sec:quickstart-remove}

When you are done with the demo, an administrator deletes the cohort \code{DEMO}:
\begin{enumerate}
\item Open the cohort \code{DEMO} in \ui{Cohorts}.
\item Click \ui{Delete\dots}. The page first says what will go: so many patients, encounters,
  bed stays, items and recordings.
\item Click \ui{Delete the demo}.
\end{enumerate}
This removes, in one step:
\begin{itemize}
\item every patient, encounter, bed stay, result, note and file link that the demo's tables
  added, and the demo's recordings, with the catalog's indexes and results of them;
\item the connection \code{demo\_omop} (the OMOP file itself is not touched);
\item the cohort \code{DEMO}. Cohorts made from it keep their members as they were, without the
  demo's patients.
\end{itemize}
Nothing else changes: what another source also holds stays, and the demo's files stay where
they are. The demo is the only cohort whose deletion also removes its patients; deleting any
other cohort leaves its patients in the catalog. Only administrators can delete the demo, and
other users see the cohort without the \ui{Delete\dots} button.

The same can be done from the command line. The first command reports what would be removed,
and the second removes it:
\begin{lstlisting}[style=shell]
medorama demo remove --dry-run
medorama demo remove
\end{lstlisting}
Every step is recorded in the audit log (\code{demo\_add}, \code{demo\_remove}).

\section{Loading the demo later}
\label{sec:quickstart-later}

\cmd{medorama demo add} loads the installer's demo into the catalog in use, for example after
\cmd{medorama init -{}-no-demo} or after deleting the demo. Only administrators may run it.
\cmd{medorama demo add DIR} loads a demo folder that you made yourself. The following commands
write the demo's files into a folder and then load them, with the demo's OMOP database and the
cohort \code{DEMO}:
\begin{lstlisting}[style=shell]
.venv/bin/python tools/make-demo-data.py ~/medorama-demo
medorama demo add ~/medorama-demo
\end{lstlisting}
A catalog holds the demo once: \cmd{demo add} refuses while the cohort \code{DEMO} exists. If the
catalog has \code{data\_roots}, the demo's folder must be inside one of them.

\file{tools/make-demo-data.py} needs NumPy and wfdb-x in \file{.venv}. Without the Python viewer
there, the following command adds them first:
\begin{lstlisting}[style=shell]
.venv/bin/pip install numpy "wfdb-x @ https://github.com/alphanumericslab/wfdb-x/archive/\
4a91916f5282fba44ca3056802fbc7def298af08.tar.gz#subdirectory=python"
\end{lstlisting}
Type it as one line, without the backslash and the line break.

Running the installer again without \opt{-{}-demo} removes \file{.venv/bin/medorama.demo}, and
catalogs created after that start empty. A catalog that already holds the demo keeps it.
\cmd{init -{}-no-demo} creates an empty catalog even while \file{medorama.demo} is there.

\section{Next}
\label{sec:quickstart-next}

Real data are databases like the synthetic ones. \Cref{part:data} explains how to import
tables, connect databases, and link recordings, ECGs and images. Running Medorama for others, as
a service, on PostgreSQL, behind HTTPS and with backups, is described in
\cref{chap:administration} and \cref{chap:deployment}. Security and privacy reviewers should read
\cref{chap:architecture}, \cref{chap:protecting} and \cref{chap:security}.
